% !TEX root = ../axiomatic.tex

\section{Related and future work}\label{s:epilogue}

\subsection{Cup-(\textit{p,i}) products}

Steenrod squares are parameterized by classes on the mod $2$ homology of $\Sym_2$.
Steenrod used this group homology viewpoint to non-constructively define -- for any prime $p$ -- operations on the mod $p$ cohomology of spaces \cite{steenrod1952reduced, steenrod1953cyclic}.
To define these constructively for odd primes, May's operadic viewpoint \cite{may1970general} was used in \cite{medina2021may_st} to provide explicit \textbf{cup-$(p,i)$ constructions}, which were implemented in the computer algebra system \href{https://github.com/ammedmar/comch}{\texttt{ComCH}} \cite{medina2021comch}.
A \mbox{cup-$(p,i)$} construction is an $\cyc_p$-equivariant chain map
\[
W(p) \to \Hom(\chains, \chains^{\ot p})
\]
inducing an isomorphism in homology where $W(p)$ is the minimal resolution of $\Fp$ as an $\Fp[\Sym_p]$-module.
Studying the moduli of cup-$(p,i)$ constructions as done here for the $p = 2$ case is left to future work.
Consult \cite{medina2024connected} for the first steps in this direction.

\subsection{Relations and secondary operations}

The cup-$i$ products provide an effective construction of coboundaries coherently witnessing the commutativity relation of the cohomological cup product at the cochain level.
These coboundaries define the squares, which are primary operations.
This comment invites the consideration of relations among the squares themselves.
There are two notable relations;
the first one, known as the \textit{Cartan relation}, expresses the interaction between these operations and the cup product:
\begin{equation*}
	\Sq^k \big( [\alpha] [\beta] \big) =
	\sum_{\mathclap{i+j=k}} \, \Sq^i\big([\alpha]\big) Sq^j\big([\beta]\big),
\end{equation*}
the second, the \textit{Adem relation} \cite{adem1952iteration}, expresses dependencies appearing through iteration:
\begin{equation*}
	\Sq^i \Sq^j =
	\sum_{k=0}^{\lfloor i/2 \rfloor} \binom{j-k-1}{i-2k} \Sq^{i+j-k} \Sq^k.
\end{equation*}
Explicit cochains witnessing them were defined in \cite{medina2020cartan} and \cite{medina2021adem} respectively, and used in the classification of low dimensional invertible fermionic topological phases of matter in \cite{kapustin2017fermionic, barkeshli2022classification}.
The analogous result for the odd prime version of the Cartan relation appears in \cite{medina2025odd_cartan}.

\subsection{Other shapes}

This paper has focused on \textit{simplicial} cup-$i$ products.
It would also be interesting to consider axiomatic characterizations of cup-$i$ products on other shapes.
Specifically, on cubes and products of general simplices.
For both of these cup-$i$ formulas already exist, see for example \cite{pilarczyk2016cubical,medina2022cube_einfty} for the cubical case, and \cite{medina2024multisimplicial} for the multisimplicial one.
The (dis)agreement of the two cubical constructions remains unsolved.

